% !TeX root = ../CV_Tramarin.tex


\sectioniten{Responsabilità scientifica per progetti di ricerca internazionali e nazionali, ammessi al finanziamento sulla base di bandi competitivi che prevedano la revisione tra pari}{Scientific Responsibility in Research Projects funded through competitive calls}

%\subsectioniten{Responsabilità scientifica in Progetti di ricerca in corso}{Ongoing Research Projects with Scientific Responsibility}

% con questo comando si sopprime la descrizione che appare sotto ogni progetto, che è invece presente nella versione in estesa del CV

\shorterCventry

\cventryit{11/2024 - 10/2025}{SILENS}
{IoT solutions for a network of navigation, logistics and “smart” vessel ergonomics services}{\newline Progetto finanziato nell’ambito del programma di ricerca dell’ecosistema dell’innovazione “I-NEST - Interconnected Nord-Est Innovation Eecosystem”, PNRR - M4C2 - Investimento 1.5 "Ecosistemi dell’innovazione per la sostenibilità”, NextGenerationEU}
{\newline
	Finanziamento complessivo: 277.783,75 €\newline
	\textbf{Ruolo: Responsabile di Unità di Ricerca} - Università di Modena e Reggio Emilia\newline
	Altri partner: Negropontelab Srl, DataVeneta Computers Srl, BeDigital Srl}{
        \textbf{Attività}: in questo progetto si intende sviluppare un sistema di monitoraggio e gestione della navigazione e della logistica
    delle imbarcazioni da diporto, in grado anche di fornire informazioni utili per la sicurezza e il comfort degli utenti a bordo. Il sistema
    prevede l’utilizzo di tecnologie IoT e di intelligenza artificiale per raccogliere e analizzare i dati provenienti da sensori e dispositivi, al fine di ottimizzare la gestione delle risorse e migliorare l'efficienza operativa. L'integrazione di queste tecnologie consente di ottenere una visione completa e in tempo reale delle operazioni, supportando decisioni informate e promuovendo la sostenibilità ambientale.}

\cventryen{11/2024 - 10/2025}{SILENS}
{IoT solutions for a network of navigation, logistics, and “smart” vessel ergonomics services}{\newline Project funded within the research program of the Innovation ecosystem “I-NEST - Interconnected Nord-Est Innovation Ecosystem,” PNRR - M4C2 - Investment 1.5 “Innovation ecosystems for sustainability,” NextGenerationEU}
{\newline
    Total funding: €277,783.75\newline
    \textbf{Role: Head of Research Unit} - University of Modena and Reggio Emilia\newline
    Other partners: Negropontelab Srl, DataVeneta Computers Srl, BeDigital Srl}{}

\cventryit{11/2023 - 02/2026}{MESMET}
{Time Sensitive Networking for future enabling technologies:
    metodi di misura e caratterizzazione metrologica in scenari ibridi
        cablati/wireless}{Progetto finanziato dal Ministero dell'Università e della Ricerca (MUR), bandi PRIN2022-PNRR (codice progetto P2022JMKEN)}
{\newline
    Finanziamento complessivo: 232.000 €\newline
	\textbf{Ruolo: Responsabile di Unità di Ricerca} - Università di Modena e Reggio Emilia\newline
	Responsabile Scientifico: Prof. Domenico Capriglione, Università di Cassino\newline
	Altri partner: Università di Brescia (prof. E. Sisinni), Consiglio Nazionale delle Ricerche CNR-IEIIT (Dott. C. Zunino)}
    {\textbf{Attività}: L'attività di ricerca prevede uno studio preliminare per aggiornare lo stato dell'arte sulle nuove applicazioni
    che richiedono TSN e sulle tecniche di misura per il contesto Time Sensitive Networking (TSN). Questa rappresenta la tecnologia abilitante per soddisfare i requisiti obbligatori dell'Industria 4.0, con la crescente adozione di sensori sofisticati e telecamere ad altissima definizione, algoritmi di apprendimento automatico e intelligenza artificiale. Inoltre, il progetto prevede la progettazione e la realizzazione di un adeguato banco di prova sperimentale e di un'ampia campagna sperimentale che consenta la progettazione e la messa a punto di metodi e procedure di misura
    per caratterizzare in modo affidabile le caratteristiche metrologiche dei sistemi coinvolti.}

\cventryen{11/2023 - 02/2026}{MESMET}
     {Time Sensitive Networking for future enabling technologies:
          measurement methods and metrological characterization in hybrid
                wired/wireless scenarios}{Project funded by the Ministry of University and Research (MUR), PRIN2022-PNRR call (project code P2022JMKEN)}
     {\newline
          Total funding: €232,000\newline
          \textbf{Role: Head of Research Unit} - University of Modena and Reggio Emilia\newline
          Principal Investigator: Prof. Domenico Capriglione, University of Cassino\newline
          Other partners: University of Brescia (Prof. E. Sisinni), National Research Council of Italy CNR-IEIIT (Dr. C. Zunino)}
          {}
    
        
\cventryit{06/2023 - 06/2025}{BRIC}{Sistema monitoraggio concentrazioni in aria e emissioni dal sottosuolo di composti volatili ID21-22}{\newline Progetto finanziato nell’ambito del Bando Ricerche in Collaborazione (BRiC -2022) dell'Istituto Nazionale per l'Assicurazione contro gli Infortuni sul Lavoro (INAIL)}
{\newline
Finanziamento complessivo: 420.000 €\newline
\textbf{Ruolo: Membro dell'Unità di Ricerca}\newline
Responsabile Scientifico: Prof. Sergio Teggi, Università di Modena e Reggio Emilia\newline
Altri Partner di progetto:\newline{}- Dipartimento di Ingegneria Civile e Ingegneria Informatica (DICII), Università
di Roma "Tor Vergata"\newline{}- Agenzia di Tutela della Salute (ATS) Città Metropolitana di Milano}{\textbf{Attività}: sviluppo e validazione di un sistema prototipale di monitoraggio in continuo delle emissioni di composti organici volatili (VOC) dal sottosuolo e delle concentrazioni di VOC nell’atmosfera immediatamente sovrastante (RT-VOC, Real-Time VOC monitoring network system) costituito da: sensori per la misura delle concentrazioni di VOC in atmosfera (CA); camere di flusso
statiche per la misura delle emissioni dal sottosuolo di VOC (FS); centralina meteo (MET);
gateway locale per la ricezione LoRaWAN dei dati CA, FS e MET e trasmissione dati 4G/5G ad un
server remoto; Dashboard per il controllo sul server remoto della rete di monitoraggio. I sensori saranno autonomi dal punto di vista dell’alimentazione elettrica ed equipaggiati per operare all’esterno. Ciò, unito alla tecnologia di trasmissione dati LoRaWAN, permetterà il posizionamento dei sensori in punti distanti anche un paio di chilometri dal Gateway locale, quindi in punti anche scarsamente accessibili, e garantirà un facile riposizionamento dei sensori.}

\cventryen{06/2023 - 06/2025}{BRIC}{Monitoring system for air concentrations and emissions from the soil of volatile compounds ID21-22}{\newline Project funded under the Collaborative Research Call (BRiC -2022) of the National Institute for Insurance against Accidents at Work (INAIL)}
{\newline
Total funding: €420,000\newline
\textbf{Role: Member of the Research Unit}\newline
Principal Investigator: Prof. Sergio Teggi, University of Modena and Reggio Emilia\newline
Other project partners:\newline{}- Department of Civil Engineering and Computer Engineering (DICII), University of Rome "Tor Vergata"\newline{}- Health Protection Agency (ATS) of the Metropolitan City of Milan}{}

%\subsectioniten{Progetti di ricerca conclusi}{Completed Research Projects}

\cventryit{2021-2023}{SHeMS}
{Smart Healthcare Modular System for Personalized Remote Health Monitoring}{Progetti finanziato tramite bando competitivo, Call no. FARMO --- 1/12/2021, Università di Modena e Reggio Emilia}
{\newline Importo totale finanziamento: 74.000 €
    \newline \textbf{Ruolo: Responsabile Scientifico}\newline
    Altri Membri: Prof.ssa Merani Maria Luisa; prof.ssa Modena Maria Grazia
    }
{\textbf{Attività}: Le attività del progetto hanno riguardato la realizzazione di una nuova generazione di Smart Healthcare Modular System (SHeMS) per il monitoraggio personalizzato della salute a distanza (RHM), sfruttando la sensor fusion, gli approcci di Machine Learning e le tecnologie Internet of Things (IoT). L'obiettivo prefissato è statao quello di assistere i pazienti (pts) e i medici nella gestione di diverse patologie, con particolare attenzione allo scompenso cardiaco (HF), e ridurre l'ospedalizzazione, registrando e fornendo un'analisi in tempo reale di parametri vitali significativi. Grazie alle caratteristiche intrinseche di comunicazione IoT, si vuole fornire una relazione continua, in tempo reale e bidirezionale tra pazienti e medici, consentendo anche la diagnosi e la riconfigurazione del sistema da remoto e la calibrazione dei sensori. Inoltre, l'inclusione dell'intelligenza artificiale (AI) all'interno del sistema consente un livello senza precedenti di funzioni di monitoraggio e di identificazione di situazioni non sicure o pericolose, attivando così in anticipo avvisi e ulteriori richieste di intervento.}

\cventryen{2021-2023}{SHeMS}
{Smart Healthcare Modular System for Personalized Remote Health Monitoring}{Project funded by a competitive call (FAR Mission Oriented, Call no. FARMO -- 1/12/2021), University of Modena and Reggio Emilia}
{\newline
    Total funding: 74.000 €\newline
    \textbf{Role: Principal investigator}\newline
    Other Members: Prof. Merani Maria Luisa; Prof. Modena Maria Grazia}{}

\cventryit{2019}{INTERACT}
{Industrial inTernet of things architEctuRes and Algorithms for time-critical Cyber-physical sysTems}{Progetti finanziato tramite bando competitivo, fondi BIRD, Università di Padova}
{\newline
    Importo totale finanziamento: 41.000 €\newline
    \textbf{Ruolo: Responsabile Scientifico}\newline
    BIRD194014 --- dal 01/05/2019 al 01/05/2021\newline
    Altri Membri: Dott. Alessandro Sona, prof. Roberto Oboe, Ing. T. Caldognetto}
{\textbf{Attività}: Le attività del progetto si inquadrano in generale nel contesto degli Smart Cyber-Physical systems per l'Industria 4.0, con particolare riferimento all'introduzione del paradigma IIoT nell'ambito dei sistemi di misura distribuita e di smart factory. La ricerca è centrata sullo sviluppo di algoritmi e strategie per la progettazione e lo sviluppo di sistemi basati su un'architettura ad intelligenza distribuita, con lo scopo di rendere l'elaborazione delle informazioni acquisite dai sensori più vicina agli stessi. Questo rende necessario ulteriori sviluppi nell'area delle reti real-time, sia cablate che wireless, per supportare lo scambio di dati di controllo e misura. Inoltre, attenzione deve essere posta nei confronti della safety, e dell'integrazione e dell'interoperabilità tra sistemi eterogenei \cite{Morato:2020a,Vitturi:2020}. La grande mole di dati raccolti da questo sistema di sensori e attuatori distribuiti richiede inoltre l'impiego di tecniche innovative di machine learning, a supporto dei classici sistemi di misura distribuita.}



\cventryen{2019}{INTERACT}
            {Industrial inTernet of things architEctuRes and Algorithms for time-critical Cyber-physical sysTems}{Project funded by a competitive call (BIRD funding, Call no. BIRD194014 --- 15/03/2019), University of Padua}
            {\newline
            Total funding: 41.000 €\newline
            \textbf{Role: Principal investigator}\newline
            Other members: Prof. Roberto Oboe, T. Caldognetto, A. Sona}{}
            

\cventryit{2017}{MAgIC}
        {Multi-Agent Intelligent Control of time-critical
        Cyber-Physical Systems over wireless}{Progetti finanziato tramite bando competitivo, fondi BIRD (BIRD175771 del 15/05/2017), Università di Padova}
        {\newline
        Importo totale finanziamento: 60.000 €\newline
            \textbf{Ruolo: Responsabile di Work Package}\newline
            Coordinatore del progetto: prof. Luca Schenato, Università di Padova\newline
            Altri membri: S. Ghidoni, S. Milani, A. Cenedese, R. Oboe}
        {\textbf{Attività}: Le attività del progetto si inquadrano in generale nel contesto del controllo di sistemi autonomi distribuiti tramite canali di comunicazione wireless con perdite \cite{Pezzutto:2018,Branz:2019a}. 
Nello specifico, il contributo personale ha riguardato l'applicazione di metodologie cross--layer,
combinando le tecniche proprie dei controlli con quelle tipiche dei sistemi di
misura distribuita ad alte prestazioni, nell'ambito di applicazioni robotiche
collaborative e automotive \cite{Branz:2019}. L'obiettivo proposto è quello di ottenere un sistema
distribuito in cui la sezione di controllo e quella di sensing e attuazione
siano in grado di scambiare dati di misura accurati a frequenze di
aggiornamento del loop di controllo fino a 1 kHz \cite{Branz:2020a}.}


\cventryen{2017}{MAgIC}
            {Multi-Agent Intelligent Control of time-critical
            Cyber-Physical Systems over wireless}{Project funded by a competitive call (BIRD funding, Call no. BIRD 175771 --- 15/05/2017), University of Padua}
            {\newline
            Total funding: 60.000 €\newline
            	\textbf{Work Package Leader}\newline
                Principal investigator: prof. Luca Schenato, University of Padova\newline
                Other members: S. Ghidoni, S. Milani, A. Cenedese, R. Oboe}
            {Responsible for the activities in two work packages of the project}

\subsectioniten{Progetti di ricerca Internazionali in fase di valutazione}{Proposals of Research Projects}

\cventryit{02/2025}{EntPowerSMG}
{Enabling Fault-Tolerant Automation Systems for Smart Microgrids by means of Time-Sensitive Networking}{\newline Proposta progettuale in fase di valutazione, a valere su bando del ``Proyectos De Generación De Conocimiento 2024'' del ``Ministerio de Ciencia, Innovación y Universidades'', Spagna}
{\newline
	Finanziamento complessivo richiesto: 228.558,00 €\newline
	\textbf{Ruolo: Membro del Gruppo di Ricerca}\newline
    Responsabile Scientifico: Prof. Manuel Alejandro Barranco González and Prof. Julián Proenza Arenas,  Universidad de las Islas Baleares, Spagna\newline
	Altri Partner di progetto:\newline{}- Università di Modena e Reggio Emilia, Italia\newline{}- Universidade do Porto, Portugal\newline{}- Mälardalen University, Sweden\newline{}- National Research Council of Italy, CNR-IEIIT\newline{}- ABB Corporate Research Center, Sweden}{}


\cventryen{02/2025}{EntPowerSMG}
    {Enabling Fault-Tolerant Automation Systems for Smart Microgrids by means of Time-Sensitive Networking}{\newline Project proposal under evaluation, submitted to the ``Proyectos De Generación De Conocimiento 2024'' call by the ``Ministerio de Ciencia, Innovación y Universidades'', Spain}
    {\newline
        Total requested funding: €228,558.00\newline
        \textbf{Role: Member of the Research Group}\newline
        Principal Investigators: Prof. Manuel Alejandro Barranco González and Prof. Julián Proenza Arenas, Universidad de las Islas Baleares, Spain\newline
        Other project partners:\newline{}- University of Modena and Reggio Emilia, Italy\newline{}- University of Porto, Portugal\newline{}- Mälardalen University, Sweden\newline{}- National Research Council of Italy, CNR-IEIIT\newline{}- ABB Corporate Research Center, Sweden}{}
    


\subsectioniten{Partecipazione in Progetti di ricerca}{Member of the research team in Research Projects}

\cventryit{2014 \& 2016}{IMET2AL}{genomIc Model prEdictive ConTrol Tools for evolutionAry pLants}{within the framework of ``Flagship project \emph{La Fabbrica del Futuro}'', CNR}{\newline
    \textbf{Ruolo: Membro del gruppo di ricerca}\newline
    Importo totale finanziamento: 312.500 €\newline
    Importo finanziamento per Unità Operativa: 104.168 €\newline
    Coordinatore del progetto: Dott. Ivan Cibrario Bertolotti, IEIIT-CNR\newline
    Lettera di conferimento incarico: Protocollo IEIIT 321 del 19/02/2014. Attività: 01/01/2014--31/12/2014, estesa a 15/01/2016 -- 15/05/2016}
{\textbf{Attività svolta:} L’attività principale all’interno del progetto di ricerca ha riguardato l'analisi e la progettazione di un’interfaccia unificata che possa essere sfruttata da un sistema di controllo industriale come punto di accesso agli strati protocollari di comunicazione, sia verso l’interno che l’esterno del sistema. In quest'ambito ci si è quindi soffermati sul problema della caratterizzazione dei ritardi introdotti dai componenti utilizzati per la realizzazione di questa interfaccia \cite{Ballarino:2014,Tramarin:2014}. Altre problematiche risiedono nella definizione di strategie di “green networking”, con la definizione di algoritmi per la predizione del traffico per \cite{Cenedese:2014}.}
%    Dr. Tramarin participated as a Member of the research Unit}


\cventryen{2014 \& 2016}{IMET2AL}{genomIc Model prEdictive conTrol Tools for evolutionAry pLants}{within the framework of ``Flagship project \emph{La Fabbrica del Futuro}'', CNR}{\newline
    Total funding: 312.500 €\newline
    Total funding per research operating Unit: 104.168 €\newline
    Principal investigator: Dott. Ivan Cibrario Bertolotti, IEIIT-CNR\newline
    Protocol no. 321 --- 19/02/2014}
{Dr. Tramarin participated as a Member of the research Unit}


\cventryen{2013-2014}{GECKO}{Generic Evolutionary Control knowledge-based Module}{within the framework of ``Flagship project  \emph{La Fabbrica del Futuro}'', CNR}{\newline
    Total funding: 700.000 €\newline
    Total funding per research operating Unit: 212.600 €\newline
    Principal investigator: Anna Valente, ITIA-CNR (then Alessandro Brusaferri, ITIA-CNR)\newline
    Protocol no. 312 --- 18/02/2014}{Dr. Tramarin participated as a Member of the research Unit}
    
\cventryit{2013-2014}{GECKO}{Generic Evolutionary Control knowledge-based Module}{within the framework of ``Flagship project  \emph{La Fabbrica del Futuro}'', CNR}{\newline
    \textbf{Ruolo: Membro del gruppo di ricerca}\newline
    Importo totale finanziamento 700.000 €\newline
    Importo finanziamento per Unità Operativa 212.600 €\newline
    Coordinatore del progetto: Anna Valente, ITIA-CNR (poi Alessandro Brusaferri, ITIA-CNR)\newline
    Lettera di conferimento incarico: Protocollo IEIIT 312 del 18/02/2014. Attività: 01/01/2013---31/12/2014}
{\textbf{Attività svolta:} all’interno del progetto di ricerca, e nello specifico dell’Unità di ricerca, l'attività svolta ha riguardato l’analisi delle problematiche di risparmio energetico, valutando le possibilità di impiego dello standard IEEE 802.3az coniugandolo con i requirements di sistemi  Real-Time basati su ethernet \cite{Tramarin:2013}. In questo ambito si è quindi presa in considerazione la  definizione di nuove strategie di risparmio energetico adatte al contesto di scambio dati real--time e proponendo nuovi metodi di implementazione dei servizi protocollari \cite{Vitturi:2015}. Si è inoltre contribuito in modo sinergico alla realizzazione degli obiettivi del progetto con gli altri membri dell’Unità operativa del CNR- IEIIT.}

\cventryit{2011-2012}{IMPROVE}{Implementing Manufacturing science solutions to increase equiPment pROductiVity and fab pErformance}{Joint Technological Initiative (JTI) ENIAC. Progetto co--finanziato MIUR (66.7\%) e Comunità
    Europea (33.3\%)}{\newline
    \textbf{Ruolo: Membro del gruppo di ricerca}\newline
    Importo totale finanziamento 18.160.970 €\newline
    Importo finanziamento per Unità Operativa 369.500 €\newline
    Coordinatore del progetto: Dott. Francois Finck, ST Microelectronics, Crolles, Francia\newline
    Responsabile CNR--IEIIT: Ing. Stefano Vitturi, IEIIT, Padova\newline
    Lettera di conferimento incarico: Protocollo IEIIT 642 data 29/04/2016. Attività: 01/01/2009---30/06/2012}
{\textbf{Attività svolta:} L’attività principale all’interno del progetto di ricerca ha riguardato la definizione di un sistema di sensori wireless, con caratteristiche di funzionamento real--time, ed  integrata con la rete di stabilimento, per il collegamento di tutti sensori e le apparecchiature ausiliarie presenti negli stabilimenti di produzione di semicondutori.
Le attività hanno inoltre riguardato l'uso di tecnologie basate su standard IEEE 802.11. Questo ha poi portato, come sviluppo, alla necessità di un’analisi delle prestazioni di sistemi di comunicazione wireless basati su tecnologia IEEE 802.11 rispetto ad indicatori tipici delle reti industriali \cite{Gamba:2009,Seno:2011,Bertocco:2012,Seno:2012}.
 I risultati ottenuti hanno anche permesso la definizione di nuove tecniche di rate adaptation per sistemi IEEE 802.11, progettate specificamente per applicazioni real--time \cite{Vitturi:2013,Tramarin:2013a}.}


\cventryen{2011-2012}{IMPROVE}{Implementing Manufacturing science solutions to increase equiPment pROductiVity and fab pErformance}{Joint Technological Initiative (JTI) ENIAC. Project co--funded by Italian Ministry of Research MIUR (66.7\%) e European Community (33.3\%)}{\newline
Total funding: 18.160.970 €\newline
Total funding per research operating Unit: 369.500 €\newline
Principal investigator: Dott. Francois Finck, ST Microelectronics, Crolles, France\newline
Protocol no. 642 --- 29/04/2016}{Dr. Tramarin participated as a Member of the research Unit}

 \cventryit{2011-2013}{PRIN 2009}{Characterisation and performance measurement in hybrid smart transducer networks: innovative experimental methods and instrumentation}{Università di Padova --- Ministero dell’Istruzione, dell’Università e della Ricerca}{\newline
\textbf{Ruolo: Membro del gruppo di ricerca}\newline
    Importo totale finanziamento 335.720 €\newline
    Importo finanziamento per Unità Operativa 81.389 € per l'Università di Padova\newline
    Importo finanziamento per Unità Operativa 78.600 € per il Consiglio Nazionale delle Ricerche\newline
    Coordinatore del progetto: prof. Claudio Narduzzi, Università di Padova\newline
    Riferimento no. 2009ZTT5N4\_001 and 2009ZTT5N4\_003 del 31/05/2010\newline Attività: 17/10/2011---17/10/2013}
{\textbf{Attività svolta:} in questo progetto di ricerca sono state valutate le prestazioni di sistemi di comunicazione wireless basati su tecnologia IEEE 802.11 rispetto ad indicatori tipici delle reti industriali \cite{Seno:2011,Seno:2012a,Bertocco:2012,Bertocco:2013b}. In quest’ambito si sono poste le basi per lo sviluppo di nuovi algoritmi di controllo automatico della velocità di trasmissione nelle reti wireless.
Inoltre, si è avviata un’importante indagine della caratterizzazione dei componenti che vengono impiegati in rete, e delle non-idealità che essi introducono, valutando ad esempio i ritardi di elaborazione \cite{Seno:2011a,Vitturi:2013a}.}

\cventryen{2011-2013}{PRIN 2009}{Characterisation and performance measurement in hybrid smart transducer networks: innovative experimental methods and instrumentation}{University of Padova --- Ministry of Instruction, University and Research}{\newline
    Total funding: 335.720 €\newline
    Total funding per research operating Unit: 81.389 € --- University of Padova\newline
    Total funding per research operating Unit: 78.600 € --- National Research Council of Italy\newline
    Principal investigator: prof. Claudio Narduzzi, University of Padova\newline
    Contract no. 2009ZTT5N4\_001 and 2009ZTT5N4\_003 --- 31/05/2010}
    {Dr. Tramarin participated as a Member of the research Unit of the University of Padova for years 2011--2012\newline
    He participated as a Member of the research Unit of the National research Council of Italy for year 2013}


    \cventryit{2010-2012}{PRIN 2008}{Measurements, and accuracy evaluation in 
    wireless space--time localization applications under real--life conditions}{Università di Padova --- Ministero dell’Istruzione, dell’Università e della Ricerca}{\newline
    \textbf{Ruolo: Membro del gruppo di ricerca}\newline
    Importo totale finanziamento 101.224 €\newline
    Importo finanziamento per Unità Operativa 24.200 € per l'Università di Padova\newline
    Coordinatore del progetto: prof. Paolo Carbone, Università di Perugia\newline
    Riferimento no. 2008TK5B55\_004 del 22/03/2010\newline Attività: 22/03/2010--22/09/2012}
{\textbf{Attività svolta:} nel contesto di questo progetto di ricerca l'attività è consistita nello sviluppo di un framework originale e completo di simulazione, per l’analisi di reti di sensori mutuamente interferenti, basate sugli standard IEEE 802.15.4 e IEEE 802.11 \cite{De-Dominicis:2009,Bertocco:2010a,Bertocco:2011a}. Questo framework è poi stato esteso e validato per comprendere sempre più scenari, in particolare verso i sistemi distribuiti real--time per applicazioni industriali.
Un altro ambito di investigazione ha riguardato la progettazione di algoritmi di ritrasmissione efficaci per reti di sensori e trasduttori wireless intelligenti \cite{Gamba:2009,Gamba:2010}.
Le attività si inquadravano nell’ambito della caratterizzazione delle capacità e dell’accuratezza nei sistemi di localizzazione mediante reti wireless di sensori, basate su misura di ToA, problematiche di sincronizzazione, in un ambiente interferente.}

\cventryen{2010-2012}{PRIN 2008}{Measurements, and accuracy evaluation in 
    wireless space--time localization applications under real--life conditions}{University of Padova --- Ministry of Instruction, University and Research}{\newline
    Total funding: 101.224 €\newline
    Total funding per research operating Unit: 24.200 € --- University of Padova\newline
    Principal investigator: prof. Paolo Carbone, University of Perugia\newline
    Contract no. 2008TK5B55\_004 --- 22/03/2010}
    {Dr. Tramarin participated as a Member of the research Unit of the University of Padova}

% si ripristina la descrizione dei progetti, che è presente nella versione estesa del CV
\longerCventry